\sectionkicker{Source-backed technical notes}
\subsection*{モデルへの入口が増える――API、open weights、独立プロバイダ: Technical Notes}
本欄は記事本文で圧縮した一次資料上の情報を、比較・再検証しやすい形で整理したものである。時系列、確認済みの主張、留意点、一次資料URLを掲載する。
\smallskip
\noindent{\small\textit{Technical Notesでは、一次資料から確認した公開・提供・研究上の事実を記載する。数値・能力評価や提供元・著者による比較表現は、独立再現済みの結果ではなく帰属claimとして扱う。}}\par
\medskip
\subsection*{Theme at a glance}
\addcontentsline{toc}{subsection}{Theme at a glance}
\begin{center}
\footnotesize
\begin{tabularx}{\linewidth}{@{}p{0.20\linewidth}p{0.10\linewidth}p{0.14\linewidth}X@{}}
\toprule
資料 & 位置づけ & 種別 & 時系列 \\
\midrule
OpenAI API no-waitlist availability & 主要資料 & API & 2021-11-18 \\
GPT-J-6B & 主要資料 & オープンウェイト & 2021-06-09 \\
Jurassic-1 / AI21 Studio & 補足資料 & API & 2021-08-04 \\
\bottomrule
\end{tabularx}
\end{center}
\normalsize

\begin{technicalnote}{OpenAI API no-waitlist availability}{主要資料}
\begingroup
% reader-facing Technical Notes break policy
\widowpenalty=10000
\clubpenalty=10000
\displaywidowpenalty=10000
\begin{tabularx}{\linewidth}{@{}>{\bfseries}p{0.22\linewidth}X@{}}
組織 & OpenAI \\
種別 & API \\
時系列 & 2021-11-18 \\
\end{tabularx}
\smallskip
{\bfseries 一次資料から整理したtechnical points}
\begin{itemize}[leftmargin=1.5em,itemsep=0.35em]
\item \textbf{一次情報で確認できる事実}: OpenAIの選定済み一次資料では、2021年のfirst-party announcementでOpenAI APIのwaitlistを撤廃し、API accessの入口を拡大したことを確認できる。 この項目はmodel weight公開ではなくhosted API availabilityの変更として扱い、open weightsやopen sourceとは区別する。
\end{itemize}
\begin{minipage}{\linewidth}
% reader-facing Technical Notes generic boundary/source tail group
\begin{samepage}
{\bfseries 一次資料}
\begingroup\sloppy
\Urlmuskip=0mu plus 2mu\relax
\begin{itemize}[leftmargin=1.5em,itemsep=0.25em]
\item {\scriptsize\url{https://openai.com/index/api-no-waitlist}}
\end{itemize}
\endgroup
\end{samepage}
\end{minipage}
\endgroup
\end{technicalnote}
\medskip

\begin{technicalnote}{GPT-J-6B}{主要資料}
\begingroup
% reader-facing Technical Notes break policy
\widowpenalty=10000
\clubpenalty=10000
\displaywidowpenalty=10000
\begin{tabularx}{\linewidth}{@{}>{\bfseries}p{0.22\linewidth}X@{}}
組織 & EleutherAI / maintainers \\
種別 & オープンウェイト \\
時系列 & 2021-06-09 \\
\end{tabularx}
\smallskip
{\bfseries 一次資料から整理したtechnical points}
\begin{itemize}[leftmargin=1.5em,itemsep=0.35em]
\item \textbf{一次情報で確認できる事実}: EleutherAI / maintainersの選定済み一次資料では、mesh-transformer-jaxの2021年commitは`gpt-j-6b release`としてGPT-J-6Bの公開境界を記録している。 本誌ではこの事実をopen-weight releaseのaccess boundaryとして扱い、hosted APIや商用serviceの提供開始とは別のevent typeに分ける。
\end{itemize}
\begin{minipage}{\linewidth}
% reader-facing Technical Notes generic boundary/source tail group
\begin{samepage}
{\bfseries 一次資料}
\begingroup\sloppy
\Urlmuskip=0mu plus 2mu\relax
\begin{itemize}[leftmargin=1.5em,itemsep=0.25em]
\item {\scriptsize\url{https://github.com/kingoflolz/mesh-transformer-jax/commit/b0c18e5205ea5ff7f97a7047e964fc04a78ea42b}}
\end{itemize}
\endgroup
\end{samepage}
\end{minipage}
\endgroup
\end{technicalnote}
\medskip

\begin{technicalnote}{Jurassic-1 / AI21 Studio}{補足資料}
\begingroup
% reader-facing Technical Notes break policy
\widowpenalty=10000
\clubpenalty=10000
\displaywidowpenalty=10000
\begin{tabularx}{\linewidth}{@{}>{\bfseries}p{0.22\linewidth}X@{}}
組織 & AI21 Labs \\
種別 & API \\
時系列 & 2021-08-04 \\
\end{tabularx}
\smallskip
{\bfseries 一次資料から整理したtechnical points}
\begin{itemize}[leftmargin=1.5em,itemsep=0.35em]
\item \textbf{一次情報で確認できる事実}: AI21 Labsの選定済み一次資料では、AI21 LabsはJurassic-1 model familyと、それへアクセスするAI21 Studioを同じfirst-party announcementで提示した。 この項目ではmodel familyそのものとprovider-operated access surfaceを分け、open-weight releaseと同一視しない。
\end{itemize}
\begin{minipage}{\linewidth}
% reader-facing Technical Notes generic boundary/source tail group
\begin{samepage}
{\bfseries 一次資料}
\begingroup\sloppy
\Urlmuskip=0mu plus 2mu\relax
\begin{itemize}[leftmargin=1.5em,itemsep=0.25em]
\item {\scriptsize\url{https://www.ai21.com/blog/announcing-ai21-studio-and-jurassic-1/}}
\end{itemize}
\endgroup
\end{samepage}
\end{minipage}
\endgroup
\end{technicalnote}
\medskip
